\documentclass[a4paper,10pt]{article}
\usepackage[margin=1.8cm,landscape]{geometry}
\usepackage[T1]{fontenc}
\usepackage{lmodern,longtable,booktabs,array}
\usepackage[hidelinks]{hyperref}
\renewcommand{\arraystretch}{1.15}
\setlength{\LTcapwidth}{\textwidth}
\renewcommand{\tablename}{Table}
\renewcommand{\thetable}{S\arabic{table}}
\begin{document}
\begin{center}
{\large\bfseries Supplementary Materials}\\[4pt]
A Measurement-Anchored Multi-Criteria Decision-Support Framework for Selecting a Blockchain Consensus Protocol\\
in IoT-Enabled Open Education\\[4pt]
\small Baikonys, Zhuzbayev, Alkhanova, Adilova, Sabitova, Tashatov and Sarsenbay
\end{center}
\noindent Both tables are generated by \texttt{make\_supplementary.py} of the repository
\url{https://github.com/magzhansaken/consensus-mcdm-open-education} (release v1.1.0): Table~S1 from
\texttt{ASSUMPTIONS.md} and Table~S2 from the calibration outputs \texttt{results/calibration\_fabric.csv} and
\texttt{results/calibrated\_params.json}.

\footnotesize
\begin{longtable}{>{\raggedright\arraybackslash}p{0.9cm} >{\raggedright\arraybackslash}p{3.2cm} >{\raggedright\arraybackslash}p{8.6cm} >{\raggedright\arraybackslash}p{2.6cm} >{\raggedright\arraybackslash}p{7.2cm}}
\caption{Assumptions of the framework. Every number that is not a published measurement comes from one of these rows.
Kind: \emph{assumption} = declared choice with its sensitivity range; \emph{documented} = specification or configuration;
\emph{calibrated} = fitted on the anchors; \emph{method} and \emph{protocol} = settings of the analysis; rows F1, A7b,
W1--W8, W-O and E1--E6 were added or extended in the revision.}\label{tab:S1}\\
\toprule ID & Item & Value & Kind & Basis\\ \midrule \endfirsthead
\toprule ID & Item & Value & Kind & Basis\\ \midrule \endhead
\bottomrule \endfoot
A1 & Consortium size & n = 7 organisations, one consensus node + one executing node each & assumption & odd n: f\_byz = 2 (BFT), f\_crash = 3 (Raft); sensitivity n = 4, 10\\
A2 & Network & campus LAN, RTT 2 ms (log-normal, CV 0.2), 1 Gbit/s links; block transfer time = bytes $\times$ fan-out / bandwidth (no per-connection window cap); sensitivity RTT 20 ms & assumption & typical campus backbone\\
A3 & Node & 4 vCPU; utilisation = CPU-seconds / (4 $\times$ wall-clock) & assumption & \\
A4 & Transaction & 1\,000 B payload + platform envelope: Fabric 2\,500 B, Besu 200 B, CometBFT 100 B, HotStuff 100 B & assumption & Fabric envelopes carry endorsements and certificates\\
A5 & Application tier & commit service c\_block + c\_tx$\cdot$B and 6.0 ms CPU per tx per executing node, calibrated on Fabric (per-role CPU of Husna et al. 2026) and applied to CometBFT and HotStuff, whose anchors have no execution layer; Besu anchors include EVM execution & calibrated + assumption & role tables of Husna et al. 2026 (peer + chaincode + CouchDB $\approx$ 6 ms CPU/tx at 50--150 TPS)\\
A6 & Block formation & Fabric: MaxMessageCount 100, BatchTimeout 2 s; Besu: block period 2 s, txpool 4\,096 tx; CometBFT: timeout\_commit 1 s, mempool 5\,000 tx (block $\leq$ 20\,000 tx); HotStuff: proposer-driven, a new proposal as soon as the 2-chain pipeline slot frees ($\geq$ 100 ms batch delay), all pending batches & documented & channel configuration of Husna et al. 2026; Besu/CometBFT docs; HotStuff benchmark settings\\
A7 & Besu capacity & QBFT, IBFT 2.0: 220 tx per 2-s block (110 TPS burst capacity); Clique: $\geq$ 423 tx per block (no backlog observed, lower bound) & calibrated & \"Ozara et al. 2026\\
A8 & Besu execution constant & o = 0.2 ms per tx (shared EVM tier) & calibrated & Clique median latency 1.0 s at a 2-s block period\\
A9 & Leader-failure handling & Raft election 5 s (TickInterval 500 ms $\times$ ElectionTick 10); SmartBFT ViewChangeTimeout 20 s; IBFT 2.0/QBFT round change 4 s (requesttimeoutseconds = 2 $\times$ blockperiodseconds, as in the Besu tutorial genesis; the documented default is 1 s and the configuration-reference example 10 s); CometBFT timeout\_propose 3 s; HotStuff timeout\_delay 5 s; Clique out-of-turn sealing 1 s & documented & Fabric sample configtx and ordering-service docs; Fabric 3 BFT configtx (ViewChangeTimeout 20s); Besu QBFT tutorial and configuration reference; CometBFT config/config.go defaults; HotStuff benchmark/fabfile.py\\
A10 & Proposer rotation & Raft stable leader; SmartBFT every 3 decisions; IBFT 2.0, QBFT, CometBFT, HotStuff, Clique every block & documented & library/protocol defaults\\
A11 & Finality & deterministic at commit for all but Clique; Clique after n/2 + 1 = 4 further blocks & documented & Clique is a probabilistic-finality PoA\\
A12 & Ledger overhead & header 400--700 B; signature 64--72 B (ECDSA/Ed25519); Fabric block metadata signature 300 B; BFT quorum certificates carry 2f + 1 signatures & documented & protocol encodings\\
A13 & Fault scenario (C7) & leader crash at 120 s and a second node crash at 300 s of a 600-s run at the deadline-day peak; no operator intervention; with n = 4 only the first crash & assumption & the injected crashes are tolerated by every protocol (n = 7 and 10: two; n = 4: one)\\
K1 & Energy proxy (C4) & 15 W per fully loaded vCPU; idle power excluded (equal for all) & assumption & \\
K2 & Cost (C11) & 0.04 USD per vCPU-hour, 0.10 USD per GB-month; nodes sized at 70 \% of peak utilisation & assumption & generic public cloud list prices\\
W1 & Active learners per weekday & 100\,000 & assumption (scale) & Open University: almost 200\,000 students; OULAD daily active share 20.9 \% (median 18.9 \%) $\rightarrow$ about 42\,000 at OU scale; base = 2.4 x OU (national multi-institution platform); sensitivity 40\,000--200\,000\\
W2 & Ledger-anchored events per active learner and day & 12 & assumption (conservative) & OULAD: 5.9 interaction records (learner x resource x day) and 22 clicks per active learner-day; 5.1 graded submissions per enrolment; sensitivity 3--24\\
W3--W4 & IoT telemetry & 5\,000 devices, one record per 120 s (41.7 tx/s) & assumption & no public campus trace; sensitivity 1\,000--20\,000 devices\\
W5 & Deadline-evening multiplier & 2.5 on 19--23 h (learner activity of the deadline day +54 \%) & assumption, supported & OULAD: interaction records on due days 1.69 x presentation mean (IQR 1.40--2.01), 2.04 x the day before; 17.6 \% of submissions on the due day; value reproducing the OULAD surge 2.94; sensitivity 1.5--4.0\\
W6 & Enrolment burst & +120 tx/s for one hour (10 h) on the stress day & assumption (stress design) & OULAD: rolling registration (99.3 \% before start, median P10--P90 span 116 days, busiest day 4.6 \%) $\rightarrow$ +0.49 tx/s at base scale; the base value stands for a centralised admission window; sensitivity 0--240 tx/s\\
W7 & Transaction size & 1\,000 B & assumption & record hash, metadata and signatures; sensitivity 250--4\,000 B\\
W8 & Hourly profile & two peaks (10--13 h, 19--23 h), busiest hour 21 h with 8.0 \% of the day; exponent 1 & assumption, supported & LeBlanc 2023 (Front. Educ. 8:1264868): peak 20--22 h, activity at every hour; Smarr \& Schirmer 2018 (Sci. Rep. 8:4793): class-day activity in the first half, sleep window 0--6 h; sensitivity exponent 0.5--1.5 (busiest hour 6.4--9.4 \%)\\
W-O & OULAD-anchored workload & W2 = 5.87, W5 = 2.94, W6 = 0.49 tx/s, other rows as above & derived from OULAD & \texttt{scripts/oulad\_workload\_stats.py} $\rightarrow$ \texttt{results/oulad\_workload\_stats.json}; levels 48.5 / 80.0 / 80.0 tx/s\\
E1--E6 & Education-facing indicators & receipt within 2 s / 5 s at the deadline peak; enrolment records per hour above the base load of the enrolment hour at 99 \% within 30 s; enrolment-hour completion; credential finality; deadline records delayed per crash incident; cost per active learner and year & definition & \texttt{consensus\_mcdm/education.py}; 2-s bound after Miller 1968\\
L1--L3 & Load levels & nominal 55.6 tx/s (daily mean), peak 108.3 tx/s (deadline-day busiest hour), stress 185 tx/s (enrolment hour) & derived from W1--W7 & \\
R1 & CometBFT transport & per-tx mempool cost o and gossip hop delay g calibrated at 175 validators (fan-out-10 gossip: 3 hops at n = 175, 1 hop at n $\leq$ 10); vote verification scales with n, mempool cost does not & calibrated + assumption & CometBFT QA v1 report\\
R2 & HotStuff anchor network & AWS multi-region WAN assumed RTT 200 ms (CV 0.2), 10 Gbit/s instances; rounds = batch dissemination + 2-chain commit (3), pipeline depth 2; proposals carry batch digests (fan-out 2) & assumption & HotStuff benchmark documentation (geo-replicated runs)\\
R3 & Calibration run protocol & Fabric fits use 90-s runs with 30-s warm-up (as in Husna et al. 2026); stage 1 fits c\_block, c\_tx, a, o, m on the batch $\times$ load campaign (n = 4); stage 2 fits $\nu$ (growth of per-tx work per added validator) on the orderer-count campaign; $\kappa$ from measured orderer CPU; 3 seeds per point for validation; RTT and crash-fault campaigns are hold-outs & protocol & \\
B1 & Besu burst plausibility check & client submission rate 150 tx/s (not reported in the anchor) & assumption & used only to report implied p50/p95\\
M1 & Weights & Entropy and CRITIC on the mean matrix; base weights = normalised geometric mean & method & project design\\
M2 & Rankings & TOPSIS (vector normalisation), VIKOR (v = 0.5), COPRAS, EDAS & method & \\
M3 & SMAA-2 & 10\,000 weight draws (uniform on the simplex; and Dirichlet centred on the base weights, concentration 50) $\times$ 30 seed matrices; aggregation by each ranking method & method & Lahdelma \& Salminen 2001\\
M4 & Seeds and CI & 30 seeds per alternative; 95 \% CI by Student t & protocol & \\
O1 & Ordinal scores C8 (attack resistance, 0--4) & one point each for: BFT safety (f < n/3), deterministic finality, proposer rotation, built-in equivocation evidence. Raft 1, Clique 1, IBFT 2.0 3, QBFT 3, SmartBFT 3, HotStuff 3, CometBFT 4 & documented & protocol specifications (Table 5 of the article)\\
O2 & Ordinal scores C10 (manageability, 0--4) & one point each for: validator-set change without restart, contract- or configuration-based validator selection, metrics endpoint, maintained production release line with load-testing tooling. Raft 4, SmartBFT 4, QBFT 4, IBFT 2.0 3 (validator change only by in-protocol voting), Clique 3 (voting only; removed in Besu 26.4), CometBFT 3 (validator set updated by the application through ABCI), HotStuff 0 (research prototype, static committee) & documented & property-level table in \texttt{consensus\_mcdm/criteria\_checks.py} (Table 5 of the article)\\
O3 & IBFT 2.0 vs QBFT & the performance anchor does not separate the two (same block period, same capacity, same round-change timer); they differ in the model only through C10 & limitation & \"Ozara et al. 2026\\
F1 & Peak-load screen & an alternative is feasible only if at least 99 \% of the stress-level (enrolment-hour) transactions commit within the 30-s client timeout in every seed; capacity = largest offered load (5-tx/s grid, 10 seeds) at which the mean committed share is $\geq$ 99 \%; screen re-checked at 95 \% and 99.9 \% & protocol (non-compensatory, conjunctive screen before the compensatory ranking) & an enrolment window is a service obligation, not a criterion to be traded off (reviewer comment); Hwang \& Yoon 1981 (conjunctive method)\\
A7b & Capacity-provisioned block formation & Besu IBFT 2.0/QBFT 440 tx per 2-s block (gas limit raised; failure timer unchanged); Fabric MaxMessageCount 300 (measured in the anchor); HotStuff proposal interval 1 s; all other settings as in A6/A9 & assumption (configuration within a published range) & Besu IBFT 2.0 reached 300--400 tx/s at a 2-s block time with a block gas limit of 672,000,000 on four validators (Hyperledger Foundation 2020); Besu QBFT/IBFT saturated at 130--150 tx/s with a 5-s block period, i.e. 650--750 tx per block (Mendes et al. 2026); Fabric batch 300 in Husna et al. 2026\\
D1 & Deployment ablations (sensitivity scenarios) & n = 4; n = 10; RTT 20 ms; no shared application tier; operator removes a crashed validator from rotation after its first detected failure (\texttt{exclude\_dead}); HotStuff with a 1-s proposal interval; Besu round-change timer 10 s instead of 4 s --- each with 10 seeds and its own Entropy $\times$ CRITIC weights & protocol & \\
\end{longtable}

\begin{longtable}{>{\raggedright\arraybackslash}p{4.4cm} >{\raggedright\arraybackslash}p{5.6cm} >{\raggedright\arraybackslash}p{1.7cm} >{\centering\arraybackslash}p{1.1cm} >{\centering\arraybackslash}p{1.7cm} >{\centering\arraybackslash}p{1.7cm} >{\raggedright\arraybackslash}p{6.0cm}}
\caption{Agreement between model and anchor campaigns: mean absolute percentage error (MAPE, \%) of throughput and
latency by set. Fit = used to estimate the constants; check and hold-out = not used in the fit. Latency is the mean
latency except for Besu (median of the burst). Simulated values are means of three seeds per point.}\label{tab:S2}\\
\toprule Platform and protocol & Set & Role & Points & MAPE throughput & MAPE latency & Note\\ \midrule \endfirsthead
\toprule Platform and protocol & Set & Role & Points & MAPE throughput & MAPE latency & Note\\ \midrule \endhead
\bottomrule \endfoot
Fabric 3.1.4, Raft & batch size $\times$ load, stable points & fit & 8 & 2.1 & 32.3 & \\
Fabric 3.1.4, Raft & batch size $\times$ load, saturated points (latency $\geq$ 5 s) & fit & 4 & 23.5 & 21.5 & measured latency depends on the run length\\
Fabric 3.1.4, Raft & number of orderers (4, 7, 10) & fit (stage 2) & 3 & 0.3 & 29.0 & \\
Fabric 3.1.4, Raft & injected round-trip time 0 ms & hold-out & 1 & 1.2 & 43.5 & measured / simulated latency 0.7\\
Fabric 3.1.4, Raft & injected round-trip time 50 ms & hold-out & 1 & 1.3 & 22.0 & measured / simulated latency 1.3\\
Fabric 3.1.4, Raft & injected round-trip time 100 ms & hold-out & 1 & 4.2 & 65.9 & measured / simulated latency 2.9\\
Fabric 3.1.4, Raft & crash campaign, no crash (baseline) & hold-out & 1 & 0.5 & 47.0 & \\
Fabric 3.1.4, Raft & crash campaign, 1-3 crashed orderers & hold-out & 3 & 0.3 & 59.6 & \\
Fabric 3.1.4, SmartBFT & batch size $\times$ load, stable points & fit & 7 & 1.4 & 21.8 & \\
Fabric 3.1.4, SmartBFT & batch size $\times$ load, saturated points (latency $\geq$ 5 s) & fit & 5 & 26.6 & 54.9 & measured latency depends on the run length\\
Fabric 3.1.4, SmartBFT & number of orderers (4, 7, 10) & fit (stage 2) & 3 & 0.6 & 34.6 & \\
Fabric 3.1.4, SmartBFT & injected round-trip time 0 ms & hold-out & 1 & 0.2 & 41.2 & measured / simulated latency 0.7\\
Fabric 3.1.4, SmartBFT & injected round-trip time 50 ms & hold-out & 1 & 2.4 & 9.4 & measured / simulated latency 1.1\\
Fabric 3.1.4, SmartBFT & injected round-trip time 100 ms & hold-out & 1 & 23.2 & 78.7 & measured / simulated latency 4.7\\
Fabric 3.1.4, SmartBFT & crash campaign, no crash (baseline) & hold-out & 1 & 0.2 & 44.7 & \\
Fabric 3.1.4, SmartBFT & crash campaign, 1-3 crashed orderers & hold-out & 3 & 18.8 & 35.7 & \\
Fabric 3.1.4, Raft and SmartBFT pooled & batch size $\times$ load, stable points & fit & 15 & 1.8 & 27.4 & values quoted in Section 4.1\\
Besu 25.7, QBFT & 5000-transaction burst & fit & 1 & 3.3 & 2.0 & latency = median; p95 error 56\%\\
Besu 25.7, IBFT 2.0 & 5000-transaction burst & fit & 1 & 3.4 & 9.4 & latency = median; p95 error 57\%\\
Besu 25.7, Clique & 5000-transaction burst & fit & 1 & -- & 11.8 & latency = median; p95 error 7\%; throughput not fitted (no backlog, capacity is a lower bound)\\
CometBFT v1 (175 validators) & offered 200 and 800 tx/s & fit & 2 & 3.7 & 12.7 & \\
CometBFT v1 (175 validators) & offered 400 tx/s & check & 1 & 8.6 & 57.6 & partially saturated run\\
CometBFT v1 (175 validators) & offered 1600 tx/s & hold-out & 1 & 8.0 & 10.1 & \\
2-chain HotStuff & 10 nodes, unsaturated (10k-100k tx/s) & fit & 3 & 3.1 & 2.7 & \\
2-chain HotStuff & 10 nodes, saturated (120k tx/s) & fit & 1 & 21.5 & 82.9 & latency not fitted\\
2-chain HotStuff & 20 and 50 nodes & hold-out & 2 & 0.2 & 43.6 & \\
2-chain HotStuff & 10 nodes, 1 and 3 crashed & hold-out & 2 & 7.4 & 63.4 & repeated leader timeouts not modelled\\
\end{longtable}
\end{document}
